\documentclass{article}

\usepackage{amsmath,amssymb}
\usepackage{xurl}
\usepackage[hidelinks]{hyperref}
\setlength{\emergencystretch}{2em}

% Shared commands for notes in docs/paper-gaps/.

\DeclareUnicodeCharacter{2080}{\ifmmode{}_{0}\else\textsubscript{0}\fi}
\DeclareUnicodeCharacter{2081}{\ifmmode{}_{1}\else\textsubscript{1}\fi}
\DeclareUnicodeCharacter{2082}{\ifmmode{}_{2}\else\textsubscript{2}\fi}
\DeclareUnicodeCharacter{2099}{\ifmmode{}_{n}\else\textsubscript{n}\fi}

\newcommand{\Lean}{\textsc{Lean}}
\ExplSyntaxOn
\NewDocumentCommand{\leanid}{m}{
  \ifmmode
    \textsf{\detokenize{#1}}
  \else
    \group_begin:
    \urlstyle{sf}
    \tl_set:Nn \l_tmpa_tl {#1}
    \tl_replace_all:Nnn \l_tmpa_tl {\_} {_}
    \exp_args:NV \path \l_tmpa_tl
    \group_end:
  \fi
}
\ExplSyntaxOff
\providecommand{\tr}{\operatorname{tr}}
\newcommand{\tnleanIssueBase}{https://github.com/LionSR/TNLean/issues/}
\newcommand{\tnleanPrBase}{https://github.com/LionSR/TNLean/pull/}
\newcommand{\tnleanDocBase}{https://sirui-lu.com/TNLean/docs/}
\newcommand{\ghissue}[1]{\href{\tnleanIssueBase#1}{GitHub issue \##1}}
\newcommand{\ghpr}[1]{\href{\tnleanPrBase#1}{PR \##1}}
\newcommand{\doclink}[1]{\href{\tnleanDocBase#1.html}{\path{#1}}}

% Dirac notation and the identity operator, as in blueprint/src/macros/common.tex.
% \providecommand keeps note-local definitions authoritative.
\usepackage{dsfont}
\providecommand{\ket}[1]{|#1\rangle}
\providecommand{\bra}[1]{\langle #1|}
\providecommand{\braket}[2]{\langle #1 | #2 \rangle}
\providecommand{\ketbra}[2]{|#1\rangle\!\langle #2|}
\providecommand{\Id}{\mathds{1}}
\providecommand{\one}{\mathds{1}}
\providecommand{\C}{\mathbb{C}}
\providecommand{\MN}[1]{M_{#1}(\C)}
\providecommand{\MD}{M_D(\C)}

% External referencing for paper sources.  The build helper writes sanitized
% auxiliary files under build/paper-aux/ and excludes duplicated source labels.
\usepackage{xr}

% Import a generated paper auxiliary file with a caller-supplied label prefix.
% The candidate paths support builds from docs/paper-gaps/, the repository root,
% and build/paper-aux/ itself.
\newcommand{\importpaperaux}[2]{%
  \IfFileExists{../../build/paper-aux/#2.aux}{%
    \externaldocument[#1]{../../build/paper-aux/#2}%
  }{%
    \IfFileExists{build/paper-aux/#2.aux}{%
      \externaldocument[#1]{build/paper-aux/#2}%
    }{%
      \IfFileExists{#2.aux}{%
        \externaldocument[#1]{#2}%
      }{}%
    }%
  }%
}

% General external reference macros
\newcommand{\paperref}[2]{\ref{#1-#2}}
\newcommand{\papereqref}[2]{(\ref{#1-#2})}

% CPSV16 source (1606.00608 / MPDO-22-12-17-2.tex) external document and shortcuts
\importpaperaux{cpsv16-}{1606.00608/MPDO-22-12-17-2-tnlean}
\newcommand{\cpsvref}[1]{\paperref{cpsv16}{#1}}
\newcommand{\cpsveqref}[1]{\papereqref{cpsv16}{#1}}

% Verdict marker: \gapnote{<kind>}{<status>}. Kind is the policy
% classification (clarification, local-correction, scope-restriction,
% unfaithful, false-source, open-gap); status is open, wip (elimination
% underway), resolved, or historical. Machine-read by the site index;
% prints a verdict line.
\newcommand{\gapnote}[2]{\par\noindent\textbf{Verdict:} #1 (#2).\par}


\title{The SAL--ZCL Blocking-Channel Conclusion of CPSV16 Proposition C.15 Is False}
\author{}
\date{2026-08-31}

\begin{document}
\maketitle
\gapnote{false-source}{resolved}

\section{Source claim}

In Appendix~C.2 of Ref.~\cite{Cirac2016MPDO_arXiv}, a tensor $K$ is fixed
which generates matrix product density operators and is both simple and in
block-injective canonical form. Its basis of normal tensors is written
$(\mathcal K_j)_j$. Zero correlation length first makes the repeated-copy
weights independent of the copy label, after which their common values are
absorbed into the representatives $\mathcal K_j$. These are the standing
hypotheses at source lines~1628--1665.

For an MPO tensor $A$ of virtual bond dimension $D$ and a virtual boundary
$X\in M_D(\mathbb C)$, write $\mathcal K_N^A(X)$ for the $N$-site physical
matrix obtained by multiplying the local matrices around the chain and
closing the virtual indices against $X$. Proposition~\cpsvref{prop2to5}
claims that saturation of the area law and zero correlation length imply the
existence of two trace-preserving completely positive maps satisfying the
diagram \cpsvref{3to5tcpm}. In particular, one of these maps must obey
\begin{align}
    \mathcal T\bigl(\mathcal K_2^K(X)\bigr)
        &= \mathcal K_3^K(X)
        \qquad\text{for every }X\in M_D(\mathbb C).
    \label{eq:cpsv16_prop_c15_blocking_channel_false_two_to_three}
\end{align}
The closing proof of Theorem~4.9 applies this map twice. Its final sentence
labels this step $(iii)\Rightarrow(v)$, whereas the theorem's displayed chain
and the hypothesis of Proposition~\cpsvref{prop2to5} place it in the logical
$(ii)\Rightarrow(v)$ route. This separate cyclic-label discrepancy is recorded
in \path{docs/paper-gaps/cpgsv17_mpdo_theorem_4_9_implication_label.tex}.
In either reading, the argument requires a trace-preserving completely
positive map satisfying
\begin{align}
    \mathcal T^{(2)}\bigl(\mathcal K_2^K(X)\bigr)
        &= \mathcal K_4^K(X)
        \qquad\text{for every }X\in M_D(\mathbb C).
    \label{eq:cpsv16_prop_c15_blocking_channel_false_two_to_four}
\end{align}
The source asserts these conclusions at lines~1810--1825 of
Ref.~\cite{Cirac2016MPDO_arXiv}.

\section{A four-letter trace-norm obstruction}

Let $l_i\in\mathbb C^2$ and $r_i\in(\mathbb C^2)^*$, for
$0\leq i<4$, be the columns and rows of
\begin{align}
    L
        &:= \begin{pmatrix}
            1&1&1&1\\
            1&2&-7&4
        \end{pmatrix},
    \label{eq:cpsv16_prop_c15_blocking_channel_false_left_pairing}\\
    R
        &:= \begin{pmatrix}
            1/4&-3/100\\
            1/4&-1/100\\
            1/4& 1/100\\
            1/4& 3/100
        \end{pmatrix}.
    \label{eq:cpsv16_prop_c15_blocking_channel_false_right_pairing}
\end{align}
Define $A_i=|l_i)(r_i|$ and let $A$ be the four-letter diagonal MPO tensor
whose only nonzero physical slices are the matrices $A_i$. The two
rectangular products satisfy
\begin{align}
    LR
        &= \begin{pmatrix}1&0\\0&0\end{pmatrix},
    \label{eq:cpsv16_prop_c15_blocking_channel_false_lr_product}\\
    T:=RL
        &= \begin{pmatrix}
            11/50&19/100&23/50&13/100\\
            6/25&23/100&8/25&21/100\\
            13/50&27/100&9/50&29/100\\
            7/25&31/100&1/25&37/100
        \end{pmatrix}.
    \label{eq:cpsv16_prop_c15_blocking_channel_false_trace_matrix}
\end{align}
Every entry of $T$ is strictly positive, and direct multiplication gives
\begin{align}
    T^2
        &= \frac14 J_4,
    \label{eq:cpsv16_prop_c15_blocking_channel_false_trace_square}\\
    T^3
        &= T^2,
    \label{eq:cpsv16_prop_c15_blocking_channel_false_trace_cube}
\end{align}
where $J_4$ is the four-by-four all-ones matrix.

Choose the virtual boundary
\begin{align}
    X
        &:= \begin{pmatrix}-1&2/5\\-5&0\end{pmatrix}.
    \label{eq:cpsv16_prop_c15_blocking_channel_false_boundary}
\end{align}
Its endpoint pairing matrix is
\begin{align}
    D:=(RXL)^{\mathsf T}
        &= \begin{pmatrix}
            0&-1/10&-1/5&-3/10\\
            1/10&0&-1/10&-1/5\\
            -4/5&-9/10&-1&-11/10\\
            3/10&1/5&1/10&0
        \end{pmatrix}.
    \label{eq:cpsv16_prop_c15_blocking_channel_false_endpoint_pairing}
\end{align}
Since $A_iA_j=T_{ij}|l_i)(r_j|$, every closure
$\mathcal K_N^A(X)$ is diagonal. Its diagonal entry at
$(i_1,\ldots,i_N)$ is
\begin{align}
    \bigl(\mathcal K_N^A(X)\bigr)_{(i_1,\ldots,i_N),(i_1,\ldots,i_N)}
        &= D_{i_1,i_N}\prod_{q=1}^{N-1}T_{i_q,i_{q+1}}.
    \label{eq:cpsv16_prop_c15_blocking_channel_false_diagonal_closure}
\end{align}
Consequently its trace norm is the sum of the absolute values of these real
diagonal entries. Using
Eqs.~\ref{eq:cpsv16_prop_c15_blocking_channel_false_trace_square} and
\ref{eq:cpsv16_prop_c15_blocking_channel_false_trace_cube}, the three
relevant values are
\begin{align}
    \left\|\mathcal K_2^A(X)\right\|_1
        &= \frac{337}{250},
    \label{eq:cpsv16_prop_c15_blocking_channel_false_two_site_norm}\\
    \left\|\mathcal K_3^A(X)\right\|_1
        &= \frac{27}{20},
    \label{eq:cpsv16_prop_c15_blocking_channel_false_three_site_norm}\\
    \left\|\mathcal K_4^A(X)\right\|_1
        &= \frac{27}{20}.
    \label{eq:cpsv16_prop_c15_blocking_channel_false_four_site_norm}
\end{align}

A trace-preserving positive map contracts the trace norm on Hermitian
matrices. The matrix $\mathcal K_2^A(X)$ is Hermitian. Thus any channel
satisfying either the two-to-three or two-to-four closure identity would give
\begin{align}
    \frac{27}{20}
        &\leq \frac{337}{250},
    \label{eq:cpsv16_prop_c15_blocking_channel_false_contractivity_contradiction}
\end{align}
which is impossible because
\begin{align}
    \frac{27}{20}-\frac{337}{250}
        &= \frac1{500}>0.
    \label{eq:cpsv16_prop_c15_blocking_channel_false_norm_gap}
\end{align}

\section{The ambient simple tensor}

It remains to place this obstruction inside the precise standing hypotheses
of Proposition~\cpsvref{prop2to5}. Let
$V:\mathbb C^4\longrightarrow\mathbb C^5$ be the coordinate isometry onto the
first four physical letters, and let $\widetilde A$ be the physical embedding
of $A$ by $V$. There are a nonzero scalar $\mu$, a normal left-canonical
bond-two tensor $B$, and an invertible bond matrix $G$ such that
\begin{align}
    0<|\mu|&<1,
    \label{eq:cpsv16_prop_c15_blocking_channel_false_weight_bound}\\
    \mu B_i
        &=G\widetilde A_iG^{-1}
        \qquad\text{for every physical pair }i.
    \label{eq:cpsv16_prop_c15_blocking_channel_false_virtual_gauge}
\end{align}
Adjoin a normal bond-one terminal tensor $C$, supported on the fifth physical
letter, and form the exact weighted direct sum
\begin{align}
    K
        &:=(\mu B)\oplus C.
    \label{eq:cpsv16_prop_c15_blocking_channel_false_ambient_tensor}
\end{align}
The two representatives $B,C$ occur once, with coefficients $\mu,1$. Their
one-site tuples span $M_2(\mathbb C)\times\mathbb C$, so this presentation is
in block-injective canonical form. Both representatives are normal, the
coefficient $1$ supplies the source normalization, and $K$ is simple. The
periodic operators generated by $K$ are positive semidefinite with strictly
positive trace at every positive length. Moreover, $K$ saturates the area law
and its physical-trace transfer is literally idempotent:
\begin{align}
    \mathcal E_K^2
        &=\mathcal E_K.
    \label{eq:cpsv16_prop_c15_blocking_channel_false_ambient_zcl}
\end{align}
Thus $K$ satisfies the complete standing simple, biCF, MPDO, SAL, ZCL, and
normalization conditions of source lines~1628--1665.

For the boundary $X$ in
Eq.~\ref{eq:cpsv16_prop_c15_blocking_channel_false_boundary}, choose the
ambient virtual boundary
\begin{align}
    Y_X
        &:=(GXG^{-1})\oplus0.
    \label{eq:cpsv16_prop_c15_blocking_channel_false_ambient_boundary}
\end{align}
The same boundary $Y_X$ is used at every length. Moving the virtual gauge to
the boundary and observing that the terminal diagonal block is zero gives
\begin{align}
    \mathcal K_N^K(Y_X)
        &=\mathcal K_N^{\widetilde A}(X)
        \qquad\text{for every }N.
    \label{eq:cpsv16_prop_c15_blocking_channel_false_boundary_isolation}
\end{align}
The physical embedding also gives
\begin{align}
    \mathcal K_N^{\widetilde A}(X)
        &=V^{\otimes N}\mathcal K_N^A(X)(V^{\otimes N})^\dagger.
    \label{eq:cpsv16_prop_c15_blocking_channel_false_physical_embedding}
\end{align}
Conjugation by $(V^{\otimes N})^\dagger$ recovers every matrix on this range.
On the orthogonal complement it can be completed by tracing and preparing a
fixed faithful density matrix. This gives a trace-preserving completely
positive recovery map on the entire five-letter matrix algebra.

Suppose now that the ambient tensor $K$ admitted the map asserted in
Eq.~\ref{eq:cpsv16_prop_c15_blocking_channel_false_two_to_three}. Apply that
map to the boundaries $Y_X$, compose on the input with the physical embedding
at length two, and compose on the output with the recovery at length three.
Equations~\ref{eq:cpsv16_prop_c15_blocking_channel_false_boundary_isolation}
and~\ref{eq:cpsv16_prop_c15_blocking_channel_false_physical_embedding} would
produce a channel from $\mathcal K_2^A(X)$ to $\mathcal K_3^A(X)$ for every
$X$. This contradicts
Eq.~\ref{eq:cpsv16_prop_c15_blocking_channel_false_contractivity_contradiction}.
The identical argument at length four rules out the twice-applied conclusion
in Eq.~\ref{eq:cpsv16_prop_c15_blocking_channel_false_two_to_four}.

The finite calculations, physical recovery, and complete standing-hypothesis
witness are formalized in the accompanying development.\footnote{The
source-facing conclusion is
\leanid{MPOTensor.CPSVBlockingChannelCounterexample.printed_propositionC15_and_theorem49_ii_to_v_are_false};
its ambient standing-hypothesis witness is
\leanid{MPOTensor.NeighboringTraceObstructionAmbientCounterexample.exists_ambient_simpleBiCF_neighboringTraceObstruction}.}

\section{Verdict}

Proposition~\cpsvref{prop2to5} is false as printed. The counterexample
satisfies every standing hypothesis used in Appendix~C.2, including the
simple biCF presentation, the MPDO condition, saturation of the area law,
literal zero correlation length, common-weight absorption, and the global
unit-weight normalization. Nevertheless, no trace-preserving completely
positive map can satisfy its two-to-three-site closure equation for every
virtual boundary.

The same witness also rules out the twice-applied two-to-four-site map used in
the logical implication (ii)$\Rightarrow$(v) of Theorem~4.9, printed with the
cyclic label (iii)$\Rightarrow$(v) in its closing proof. Conditional
channel constructions from supplied neighboring-trace factorizations remain
valid, but those extra factorizations are not consequences of SAL and ZCL.

\bibliographystyle{alpha}
\bibliography{references}

\end{document}
